\slajdid{09_3-s068}
\begin{frame}[label=09_3-s068]{Dinamička logika}
\gusttekst\setlength{\abovedisplayskip}{3pt}\setlength{\belowdisplayskip}{3pt}\setlength{\abovedisplayshortskip}{2pt}\setlength{\belowdisplayshortskip}{2pt}Problem koji se pojavljuje kod složenih CMOS logičkih kola jeste velik broj pMOS tranzistora sa velikim površinama u PUN mreži.
\par\medskip\textcolor{DEcrvena}{To može da se reši pseudo NMOS tehnologijom ali na taj način dobijamo veliku disipaciju kada je logička nula na izlazu.}
\par\medskip
U situaciji kada raspolažemo sinhronizacionim signalom sa kojim treba da radi naš digitalni sistem jedno od mogućih rešenja jeste i dinamička logika
\begin{columns}[c,onlytextwidth]\column{.48\textwidth}\centering\input{slike/tikz/s068_dinamicka_mreza.tex}\column{.5\textwidth}\centering\begin{tikzpicture}[x=.65cm,y=.65cm,font=\sitneoznake,>=Latex]\draw[->](0,2)--(6.3,2)node[below]{t};\draw[->](.2,1.8)--(.2,3.6)node[left]{CLK};\draw[->](0,0)--(6.3,0)node[below]{t};\draw[->](.2,-.2)--(.2,1.6)node[left]{Y};\draw(.2,2)--(.6,2)--(.6,3.2)--(1.6,3.2)--(1.6,2)--(2.6,2)--(2.6,3.2)--(3.6,3.2)--(3.6,2)--(4.6,2)--(4.6,3.2)--(5.6,3.2)--(5.6,2);\foreach\x in{.6,1.6,2.6,3.6,4.6,5.6}{\draw[dashed](\x,2)--(\x,-.15);}\node[below]at(1.1,0){Y=0};\node[below]at(3.1,0){Y=1};\node[below]at(5.1,0){Y=0};\draw(.2,1.2)--(.6,1.2)--(1.2,0)--(1.6,0)--(2.2,1.2)--(4.6,1.2)--(5.2,0)--(5.6,0);\end{tikzpicture}
\end{columns}
\end{frame}
